\documentclass[11pt,reqno]{amsart}

\usepackage[T1]{fontenc}
\usepackage{amsmath,amssymb,amsthm}
\usepackage[margin=1.15in]{geometry}
\usepackage{booktabs}
\usepackage{microtype}
\usepackage[colorlinks=true,linkcolor=black,citecolor=black,urlcolor=blue]{hyperref}

\theoremstyle{plain}
\newtheorem{theorem}{Theorem}[section]
\newtheorem{proposition}[theorem]{Proposition}
\newtheorem{lemma}[theorem]{Lemma}
\newtheorem{corollary}[theorem]{Corollary}
\newtheorem{conjecture}[theorem]{Conjecture}
\theoremstyle{definition}
\newtheorem{definition}[theorem]{Definition}
\newtheorem{remark}[theorem]{Remark}
\newtheorem{observation}[theorem]{Observation}
\newtheorem{example}[theorem]{Example}
\newtheorem{question}[theorem]{Question}

\newcommand{\Z}{\mathbb{Z}}
\newcommand{\N}{\mathbb{N}}
\newcommand{\Q}{\mathbb{Q}}
\newcommand{\Uu}{\mathcal{U}}
\newcommand{\Tr}{\operatorname{Tr}}
\newcommand{\Ker}{\operatorname{Ker}}

\begin{document}

\title[Knuth products in linear numeration systems]
      {Knuth products in linear numeration systems:\\
       associativity is closure}

\author{}
\date{13 August 2026}

\begin{abstract}
Knuth's \emph{Fibonacci multiplication} $n\circ m=\sum_{i,j}F_{k_i+l_j}$, defined on Zeckendorf
expansions $n=\sum_iF_{k_i}$ and $m=\sum_jF_{l_j}$, is associative. We show that this is an instance
of a single exact identity valid in every linear numeration system, under no hypotheses beyond the
defining recurrence being the minimal one of its root: if $\Phi$ reads greedy expansions into
$\Z[\beta]$ and $L\colon\Z[\beta]\to\Z$ is the unique $\Z$-linear map with $L(\beta^{k})=u_{k}$, then
the shifted Knuth product satisfies
\[
  n\circ_{s}m \;=\; L\bigl(\beta^{s}\,\Phi(n)\,\Phi(m)\bigr).
\]
Consequently $\circ_{s}$ is associative as soon as the set $S=\Phi(\N)$ of ``$\Uu$-integers''
satisfies the closure condition $\beta^{s}S\cdot S\subseteq S$, in which case $(\N,\circ_{s})$ is
isomorphic to a multiplicative subsemigroup of $\Z[\beta]$; and in general the failure of
associativity is measured by an exact bilinear obstruction. Earlier work treats numeration systems
with canonical initial values, for which the set of admissible shifts is necessarily an interval. We
drop that restriction, and the resulting extra freedom --- a free parameter $\lambda\in\Q(\beta)$,
since $L=\Tr(\lambda\,\cdot\,)$ --- produces behaviour with no canonical counterpart. For the Lucas
system $1,3,4,7,11,\dots$ the product $\circ_{s}$ is associative for $s=2$, \emph{fails} for $s=3$,
and is associative again for $4\le s\le14$: the set of admissible shifts need not be upward closed.
It also depends on the initial values and not only on $\beta$: the system $1,2,5,12,29,\dots$
admits $s=1$, whereas the canonical system for the same $\beta=1+\sqrt2$ does not. All assertions
of a finite nature are certified by exact computation; the certificates are reproduced in the text.
\end{abstract}

\maketitle

\section{Introduction}

Let $F_{1}=F_{2}=1$, $F_{k+1}=F_{k}+F_{k-1}$. Every $n\in\N$ has a unique Zeckendorf expansion
$n=\sum_{i}F_{k_{i}}$ with $k_{i}\ge2$ and $k_{i+1}\ge k_{i}+2$ \cite{Zeckendorf}. Knuth
\cite{Knuth} defined
\begin{equation}\label{eq:knuth}
  n\circ m \;=\; \sum_{i,j}F_{k_{i}+l_{j}},
  \qquad n=\sum_{i}F_{k_{i}},\quad m=\sum_{j}F_{l_{j}},
\end{equation}
and proved the unexpected fact that $\circ$ is associative. The operation is visibly commutative,
and there is no reason of any obvious kind why regrouping should be harmless: the multiset
$\{k_{i}+l_{j}\}$ is in general not a Zeckendorf expansion, so the right-hand side of
\eqref{eq:knuth} must be renormalised before the next product is formed, and renormalisation is not
a linear operation.

Arnoux \cite{Arnoux} explained the phenomenon for the Fibonacci case by transporting the
multiplication of $\Z[\varphi]$, $\varphi=(1+\sqrt5)/2$, through the Zeckendorf expansion. The
construction was subsequently extended to linear recurrences with \emph{canonical} initial values
\cite{Grabner}, where associativity is obtained once an index shift is chosen large enough, and to
Pisot numeration systems by Messaoudi \cite{Messaoudi1,Messaoudi2}. In the canonical case the
greedy expansions are exactly the $\beta$-admissible digit strings \cite{Bertrand}, the image of the
expansion map is the set $\Z_{\beta}^{\ge0}$ of non-negative $\beta$-integers, and the least
admissible shift is the constant $L^{\otimes}(\beta)$ measuring how many fractional digits a product
of two $\beta$-integers can require \cite{BFGK,AFMP}.

The purpose of this note is to isolate what is actually responsible for associativity, and to record
what happens when the classical hypotheses are removed. Our standing assumptions are minimal:
$\beta$ is an algebraic integer with monic minimal polynomial
$\chi(x)=x^{d}-a_{1}x^{d-1}-\dots-a_{d}\in\Z[x]$, and $\Uu=(u_{k})_{k\ge0}$ is any strictly
increasing sequence of positive integers with $u_{0}=1$ obeying the associated recurrence. Nothing
is assumed about $\beta$ being Pisot or Parry, nothing about admissibility of digit strings, and ---
crucially --- nothing about the initial values $u_{0},\dots,u_{d-1}$ beyond monotonicity.

Theorem \ref{thm:transport} states the identity
$n\circ_{s}m=L(\beta^{s}\Phi(n)\Phi(m))$ referred to in the abstract; it is an exercise, but it is
the whole mechanism. Theorem \ref{thm:criterion} converts it into a criterion: writing
$S=\Phi(\N)$, the closure condition $\beta^{s}S\cdot S\subseteq S$ is equivalent to the vanishing of
the \emph{defect} $\delta_{s}(n,m)=\Phi(n\circ_{s}m)-\beta^{s}\Phi(n)\Phi(m)$, and implies that
$n\mapsto\beta^{s}\Phi(n)$ is an isomorphism of $(\N,\circ_{s})$ onto a multiplicative subsemigroup
of $\Z[\beta]$. Part (iv) of that theorem is an exact identity: associativity holds if and only if
\[
  L\bigl(\beta^{s}\delta_{s}(n,m)\Phi(p)\bigr)=L\bigl(\beta^{s}\Phi(n)\delta_{s}(m,p)\bigr)
  \qquad\text{for all }n,m,p .
\]

Section \ref{sec:phenomena} contains the point of the exercise. Lemma \ref{lem:multiples} shows that
the set of shifts satisfying the closure condition is closed under multiplication by positive
integers, and Remark \ref{rem:canonical} shows that in the classical (canonical) situation one has
$\beta S\subseteq S$, so that the admissible shifts form an interval $[s_{0},\infty)$ and no finer
structure can appear. Without canonicity this collapses:

\begin{itemize}
\item For the Lucas system $u=(1,3,4,7,11,18,\dots)$ the product $\circ_{2}$ is associative,
      $\circ_{3}$ is \emph{not}, and $\circ_{s}$ is associative again for all $4\le s\le14$
      (Proposition \ref{prop:lucas} gives the certificate
      $(1\circ_{3}2)\circ_{3}6=268\ne278=1\circ_{3}(2\circ_{3}6)$).
\item For $u=(1,4,5,9,14,23,\dots)$ --- again $\beta=\varphi$ --- the admissible shifts observed are
      $\{4\}\cup[6,14]$, with $5$ excluded.
\item The set of admissible shifts is not a function of $\beta$: three systems with $\beta=\varphi$
      give $[2,\infty)$, $\{2\}\cup[4,\infty)$ and $\{4\}\cup[6,\infty)$ respectively (as far as
      computed), and for $\beta=1+\sqrt2$ the non-canonical system $1,2,5,12,29,\dots$ admits the
      shift $s=1$ while the canonical system $1,3,7,17,41,\dots$ does not.
\end{itemize}

Finally, in every system and every shift examined, associativity and closure were found to be
equivalent. We have not been able to prove the missing implication, and state it as Conjecture
\ref{conj:noaccident}: an associative Knuth product is always a transported ring product.

All finite assertions in this paper --- in particular every negative result --- are exact integer
computations, and are reproduced by the script \texttt{verify.py} accompanying this note; see
Appendix \ref{sec:verification}.

\section{Setting}

\begin{definition}\label{def:system}
A \emph{linear numeration system} is a pair $(\beta,\Uu)$ subject to:
\begin{enumerate}
\item[(H1)] $\beta$ is an algebraic integer whose minimal polynomial over $\Q$ is
      $\chi(x)=x^{d}-a_{1}x^{d-1}-\dots-a_{d}$ with $a_{i}\in\Z$;
\item[(H2)] $\Uu=(u_{k})_{k\ge0}$ is a sequence of integers with $u_{0}=1$,
      $u_{0}<u_{1}<u_{2}<\dots$, and
      $u_{k}=a_{1}u_{k-1}+a_{2}u_{k-2}+\dots+a_{d}u_{k-d}$ for all $k\ge d$.
\end{enumerate}
\end{definition}

The initial values $u_{0}=1,u_{1},\dots,u_{d-1}$ are otherwise unconstrained. Taking
$u_{k}=a_{1}u_{k-1}+\dots+a_{k}u_{0}+1$ for $1\le k<d$ gives the \emph{canonical} system associated
with the digit string $a_{1}a_{2}\cdots a_{d}$; this is the case treated in
\cite{Grabner,Messaoudi1,Messaoudi2,AFMP}, and it is not the case of interest here.

\begin{definition}\label{def:greedy}
The \emph{greedy expansion} $K(n)$ of $n\in\N$ is the finite multiset of non-negative integers
defined by $K(0)=\varnothing$ and, for $n\ge1$,
\[
  K(n)=\{\kappa(n)\}\uplus K\bigl(n-u_{\kappa(n)}\bigr),
  \qquad \kappa(n)=\max\{k:u_{k}\le n\}.
\]
\end{definition}

Since $u_{0}=1$ the set $\{k:u_{k}\le n\}$ is non-empty, and it is finite because $(u_{k})$ is
strictly increasing; as $n-u_{\kappa(n)}<n$ the recursion terminates. Thus $K(n)$ is well defined,
finite, and $n=\sum_{k\in K(n)}u_{k}$, the sum being taken with multiplicity. Multiplicities do
occur: in the Lucas system $1,3,4,7,\dots$ one has $K(2)=\{0,0\}$.

\begin{definition}\label{def:product}
For $s\ge0$ the \emph{Knuth product of shift $s$} is
\[
  n\circ_{s}m \;=\; \sum_{k\in K(n)}\ \sum_{l\in K(m)}u_{k+l+s}
  \qquad (n,m\in\N),
\]
sums taken with multiplicity. We write $\Phi(n)=\sum_{k\in K(n)}\beta^{k}\in\Z[\beta]$ and
$S=\Phi(\N)$.
\end{definition}

The operation $\circ_{s}$ is commutative, and $n\circ_{s}0=0$.

\begin{example}\label{ex:knuth}
For $u_{k}=F_{k+2}$, hypothesis (H2) holds with $\chi(x)=x^{2}-x-1$, $\beta=\varphi$, and $K(n)$ is
the Zeckendorf expansion of $n$ with indices shifted down by $2$. If $n=\sum_{i}F_{K_{i}}$ and
$m=\sum_{j}F_{L_{j}}$ in Knuth's indexing, then $K_{i}=k_{i}+2$ and $L_{j}=l_{j}+2$, so
$F_{K_{i}+L_{j}}=F_{k_{i}+l_{j}+4}=u_{k_{i}+l_{j}+2}$. Hence Knuth's operation \eqref{eq:knuth} is
exactly $\circ_{2}$ for this system.
\end{example}

\begin{example}\label{ex:base}
For $u_{k}=b^{k}$ ($b\ge2$) we have $d=1$, $\chi(x)=x-b$, $\beta=b$, and $K(n)$ lists the positions
of the base-$b$ digits of $n$ with multiplicity. Then $n\circ_{s}m=b^{s}nm$.
\end{example}

\section{The transport identity}

\begin{theorem}\label{thm:transport}
Let $(\beta,\Uu)$ be a linear numeration system. There is a unique $\Z$-linear map
$L\colon\Z[\beta]\to\Z$ with $L(\beta^{k})=u_{k}$ for all $k\ge0$, and it satisfies
\begin{enumerate}
\item[(i)] $L(\Phi(n))=n$ for every $n\in\N$;
\item[(ii)] $n\circ_{s}m=L\bigl(\beta^{s}\Phi(n)\Phi(m)\bigr)$ for all $n,m\in\N$ and all $s\ge0$.
\end{enumerate}
\end{theorem}

\begin{proof}
Since $\chi$ is monic, division with remainder by $\chi$ is available in $\Z[x]$; and since $\chi$ is
the minimal polynomial of $\beta$, a polynomial $P\in\Z[x]$ satisfies $P(\beta)=0$ if and only if
$\chi\mid P$ in $\Q[x]$, hence (writing $P=Q\chi+R$ in $\Z[x]$ with $\deg R<d$ and noting $R(\beta)=0$
forces $R=0$) if and only if $\chi\mid P$ in $\Z[x]$. Therefore the evaluation map induces an
isomorphism $\Z[x]/(\chi)\xrightarrow{\ \sim\ }\Z[\beta]$, and $\Z[\beta]$ is a free $\Z$-module with
basis $1,\beta,\dots,\beta^{d-1}$.

Define $L$ on this basis by $L(\beta^{i})=u_{i}$ for $0\le i<d$ and extend $\Z$-linearly; any map as
in the statement must agree with this one, whence uniqueness. We check $L(\beta^{k})=u_{k}$ for
$k\ge d$ by strong induction. Multiplying $\chi(\beta)=0$ by $\beta^{k-d}$ gives
$\beta^{k}=a_{1}\beta^{k-1}+\dots+a_{d}\beta^{k-d}$, and all exponents on the right lie in
$\{k-d,\dots,k-1\}\subseteq\Z_{\ge0}$, so by linearity and the inductive hypothesis
\[
  L(\beta^{k})=a_{1}u_{k-1}+\dots+a_{d}u_{k-d}=u_{k},
\]
the last equality by (H2).

(i) is immediate: $L(\Phi(n))=\sum_{k\in K(n)}L(\beta^{k})=\sum_{k\in K(n)}u_{k}=n$. For (ii),
\[
  \beta^{s}\Phi(n)\Phi(m)=\sum_{k\in K(n)}\sum_{l\in K(m)}\beta^{k+l+s},
\]
and applying $L$ termwise gives $\sum_{k,l}u_{k+l+s}=n\circ_{s}m$.
\end{proof}

\begin{corollary}\label{cor:bijection}
$\Phi\colon\N\to S$ is a bijection, with inverse $L|_{S}$. Moreover $\Phi(0)=0$ and $\Phi(1)=1$.
\end{corollary}

\begin{proof}
Surjectivity onto $S=\Phi(\N)$ holds by definition, and injectivity follows from
Theorem \ref{thm:transport}(i), which also identifies the inverse. As $u_{1}>u_{0}=1$ we have
$\kappa(1)=0$, so $K(1)=\{0\}$ and $\Phi(1)=\beta^{0}=1$.
\end{proof}

We record two structural remarks; neither is used later, but both clarify what the data of a
numeration system amounts to.

\begin{proposition}\label{prop:trace}
There is a unique $\lambda\in\Q(\beta)$ with
$L(z)=\Tr_{\Q(\beta)/\Q}(\lambda z)$ for all $z\in\Z[\beta]$. For $u_{k}=F_{k+2}$ one has
$\lambda=\varphi^{2}/\sqrt5$, and for $u_{k}=L_{k+1}$ (the Lucas system $1,3,4,7,\dots$) one has
$\lambda=\varphi$.
\end{proposition}

\begin{proof}
$L$ extends uniquely to a $\Q$-linear functional on $\Q(\beta)=\Z[\beta]\otimes_{\Z}\Q$. As
$\Q(\beta)/\Q$ is separable, the trace pairing $(x,y)\mapsto\Tr(xy)$ is non-degenerate, so
$\lambda\mapsto\Tr(\lambda\,\cdot\,)$ is an isomorphism from $\Q(\beta)$ onto its $\Q$-dual; take
$\lambda$ to be the preimage of the extension of $L$. For the two examples, with
$\psi=(1-\sqrt5)/2$ one has $\Tr(\varphi\cdot\varphi^{k})=\varphi^{k+1}+\psi^{k+1}=L_{k+1}$ and
$\Tr\bigl((\varphi^{2}/\sqrt5)\varphi^{k}\bigr)=(\varphi^{k+2}-\psi^{k+2})/\sqrt5=F_{k+2}$.
\end{proof}

Thus, for fixed $\beta$, choosing the initial values of $\Uu$ is choosing $\lambda$, and the
classical theory is the study of one distinguished $\lambda$.

\begin{remark}\label{rem:galois}
The construction is equivariant for the Galois action. If $\sigma$ is an embedding of $\Q(\beta)$
into $\overline{\Q}$ and $\beta'=\sigma(\beta)$, then $(\beta',\Uu)$ is again a linear numeration
system, $\Phi_{\beta'}=\sigma\circ\Phi_{\beta}$, $L_{\beta'}\circ\sigma=L_{\beta}$, and
$\sigma(S_{\beta})=S_{\beta'}$. In particular the closure condition of Theorem \ref{thm:criterion}
holds for $\beta$ if and only if it holds for any conjugate, so the choice of root of $\chi$ is
immaterial. Note also that $\circ_{s}$ itself never depended on that choice.
\end{remark}

\section{Associativity is closure}

\begin{definition}\label{def:defect}
For $s\ge0$ the \emph{defect} of $\circ_{s}$ is
\[
  \delta_{s}(n,m)\;=\;\Phi(n\circ_{s}m)-\beta^{s}\Phi(n)\Phi(m)\;\in\;\Z[\beta].
\]
\end{definition}

\begin{theorem}\label{thm:criterion}
Let $(\beta,\Uu)$ be a linear numeration system and $s\ge0$. Then:
\begin{enumerate}
\item[(i)] $\delta_{s}(n,m)\in\Ker L$ for all $n,m\in\N$;
\item[(ii)] $\beta^{s}S\cdot S\subseteq S$ if and only if $\delta_{s}\equiv0$ on $\N\times\N$;
\item[(iii)] if $\beta^{s}S\cdot S\subseteq S$, then $\Psi:=\beta^{s}\Phi$ is injective and satisfies
      $\Psi(n\circ_{s}m)=\Psi(n)\Psi(m)$; consequently $(\N,\circ_{s})$ is isomorphic, as a
      commutative semigroup, to $\bigl(\Psi(\N),\times\bigr)\subseteq(\Z[\beta],\times)$, and
      $\circ_{s}$ is associative;
\item[(iv)] $\circ_{s}$ is associative if and only if
      \begin{equation}\label{eq:obstruction}
        L\bigl(\beta^{s}\,\delta_{s}(n,m)\,\Phi(p)\bigr)
        =L\bigl(\beta^{s}\,\Phi(n)\,\delta_{s}(m,p)\bigr)
        \qquad\text{for all }n,m,p\in\N .
      \end{equation}
\end{enumerate}
\end{theorem}

\begin{proof}
(i) By Theorem \ref{thm:transport}(i), $L(\Phi(n\circ_{s}m))=n\circ_{s}m$, and by
Theorem \ref{thm:transport}(ii), $L(\beta^{s}\Phi(n)\Phi(m))=n\circ_{s}m$ as well; subtract.

(ii) Assume $\beta^{s}S\cdot S\subseteq S$ and put $z=\beta^{s}\Phi(n)\Phi(m)$. Then $z\in S$, say
$z=\Phi(N)$ with $N\in\N$, and Corollary \ref{cor:bijection} gives $N=L(z)$, which equals
$n\circ_{s}m$ by Theorem \ref{thm:transport}(ii). Hence $\Phi(n\circ_{s}m)=\Phi(N)=z$, i.e.
$\delta_{s}(n,m)=0$. Conversely, if $\delta_{s}\equiv0$ then
$\beta^{s}\Phi(n)\Phi(m)=\Phi(n\circ_{s}m)\in S$ for all $n,m$, which is the stated inclusion.

(iii) By (ii), $\Phi(n\circ_{s}m)=\beta^{s}\Phi(n)\Phi(m)$, so
\[
  \Psi(n\circ_{s}m)=\beta^{s}\Phi(n\circ_{s}m)=\beta^{s}\cdot\beta^{s}\Phi(n)\Phi(m)=\Psi(n)\Psi(m).
\]
$\Phi$ is injective by Corollary \ref{cor:bijection} and $\beta\ne0$ is a non-zero-divisor of the
integral domain $\Z[\beta]$, so $\Psi$ is injective. Therefore
\[
  \Psi\bigl((n\circ_{s}m)\circ_{s}p\bigr)=\Psi(n)\Psi(m)\Psi(p)
  =\Psi\bigl(n\circ_{s}(m\circ_{s}p)\bigr),
\]
and injectivity of $\Psi$ yields associativity. That $\Psi$ is a semigroup isomorphism onto its
image is the displayed multiplicativity together with injectivity.

(iv) By Theorem \ref{thm:transport}(ii) applied to the pair $(n\circ_{s}m,\,p)$ and
Definition \ref{def:defect},
\begin{align*}
  (n\circ_{s}m)\circ_{s}p
  &=L\bigl(\beta^{s}\Phi(n\circ_{s}m)\Phi(p)\bigr)\\
  &=L\bigl(\beta^{s}\bigl[\beta^{s}\Phi(n)\Phi(m)+\delta_{s}(n,m)\bigr]\Phi(p)\bigr)\\
  &=L\bigl(\beta^{2s}\Phi(n)\Phi(m)\Phi(p)\bigr)+L\bigl(\beta^{s}\delta_{s}(n,m)\Phi(p)\bigr),
\end{align*}
and symmetrically
\[
  n\circ_{s}(m\circ_{s}p)
  =L\bigl(\beta^{2s}\Phi(n)\Phi(m)\Phi(p)\bigr)+L\bigl(\beta^{s}\Phi(n)\delta_{s}(m,p)\bigr).
\]
Subtracting gives the claim.
\end{proof}

\begin{definition}\label{def:sets}
Put
\[
  C(\Uu)=\{s\ge0:\beta^{s}S\cdot S\subseteq S\},\qquad
  G(\Uu)=\{s\ge0:\circ_{s}\text{ is associative}\},
\]
so that $C(\Uu)\subseteq G(\Uu)$ by Theorem \ref{thm:criterion}(iii).
\end{definition}

\begin{lemma}\label{lem:multiples}
If $s\in C(\Uu)$ then $\beta^{s}S\subseteq S$, and $ks\in C(\Uu)$ for every integer $k\ge1$. Thus
$C(\Uu)$ is closed under multiplication by positive integers.
\end{lemma}

\begin{proof}
By Corollary \ref{cor:bijection}, $1=\Phi(1)\in S$, so
$\beta^{s}S=\beta^{s}S\cdot\{1\}\subseteq\beta^{s}S\cdot S\subseteq S$. Iterating,
$\beta^{js}S\subseteq S$ for all $j\ge0$, whence
\[
  \beta^{(j+1)s}S\cdot S=\beta^{js}\bigl(\beta^{s}S\cdot S\bigr)\subseteq\beta^{js}S\subseteq S. \qedhere
\]
\end{proof}

\begin{remark}\label{rem:canonical}
Suppose $\Uu$ is the canonical system of a simple Parry number $\beta$ \cite{Parry}. Then greedy
$\Uu$-expansions coincide with the $\beta$-admissible digit strings \cite{Bertrand}, so that
$S=\Z_{\beta}^{\ge0}$; multiplying an admissible string by $\beta$ appends a zero digit at the
bottom and preserves admissibility, whence $\beta S\subseteq S$. Combining this with
Lemma \ref{lem:multiples}, $C(\Uu)$ is upward closed, i.e.\ an interval $[s_{0},\infty)$, and
$s_{0}=L^{\otimes}(\beta)$ in the notation of \cite{BFGK,AFMP}. Consequently none of the phenomena
of Section \ref{sec:phenomena} can be seen in the canonical setting: there, ``$s$ large enough'' is
the whole story, which is the form the result takes in \cite{Grabner,Messaoudi1,Messaoudi2}.
\end{remark}

\begin{example}\label{ex:knuththeorem}
Knuth's theorem is the assertion $2\in C(\Uu)$ for the Zeckendorf system, i.e.
$\varphi^{2}\,\Z_{\varphi}^{\ge0}\!\cdot\Z_{\varphi}^{\ge0}\subseteq\Z_{\varphi}^{\ge0}$: a product
of two non-negative golden-mean integers needs at most two golden-mean fractional digits. By
Remark \ref{rem:canonical}, $C=[2,\infty)$ for that system, and Theorem \ref{thm:criterion}(iii)
identifies $(\N,\circ_{2})$ with the semigroup
$\varphi^{2}\Z_{\varphi}^{\ge0}$ under multiplication.
\end{example}

\section{Phenomena outside the canonical case}\label{sec:phenomena}

For $N\ge1$ let $G_{N}(\Uu)$ denote the set of $s$ for which associativity holds for all
$n,m,p\le N$, and $C_{N}(\Uu)$ the set of $s$ for which $\beta^{s}\Phi(n)\Phi(m)\in S$ for all
$n,m\le N$. Thus $G(\Uu)\subseteq G_{N}(\Uu)$ and $C(\Uu)\subseteq C_{N}(\Uu)$, and an exhibited
$s\notin G_{N}(\Uu)$ \emph{proves} $s\notin G(\Uu)$.

Table \ref{tab:main} lists $C_{24}\cap[0,8]$ and $G_{24}\cap[0,8]$ for thirteen systems; the two
coincided in every instance. Table \ref{tab:focus} extends the range for five systems.

\begin{table}[ht]
\centering
\small
\begin{tabular}{llllc}
\toprule
system & $u_{0},u_{1},u_{2},u_{3},\dots$ & $\chi$ & $\beta$ & $C_{24}=G_{24}$ in $[0,8]$\\
\midrule
base $2$        & $1,2,4,8,\dots$      & $x-2$        & $2$      & $0,1,2,\dots,8$\\
base $3$        & $1,3,9,27,\dots$     & $x-3$        & $3$      & $0,1,2,\dots,8$\\
Zeckendorf      & $1,2,3,5,8,\dots$    & $x^{2}-x-1$  & $1.6180$ & $2,3,\dots,8$\\
Lucas           & $1,3,4,7,11,\dots$   & $x^{2}-x-1$  & $1.6180$ & $2,\ \ 4,5,\dots,8$\\
shifted Fib.    & $1,4,5,9,14,\dots$   & $x^{2}-x-1$  & $1.6180$ & $4,\ \ 6,7,8$\\
Pell canonical  & $1,3,7,17,41,\dots$  & $x^{2}-2x-1$ & $2.4142$ & $2,3,\dots,8$\\
Pell natural    & $1,2,5,12,29,\dots$  & $x^{2}-2x-1$ & $2.4142$ & $1,2,\dots,8$\\
$x^{2}=3x+1$    & $1,4,13,43,\dots$    & $x^{2}-3x-1$ & $3.3028$ & $2,3,\dots,8$\\
$x^{2}=2x+2$    & $1,3,8,22,\dots$     & $x^{2}-2x-2$ & $2.7321$ & $4,5,\dots,8$\\
Tribonacci      & $1,2,4,7,13,\dots$   & $x^{3}-x^{2}-x-1$ & $1.8393$ & $3,4,\dots,8$\\
$x^{3}=x^{2}+1$ & $1,2,3,4,6,9,\dots$  & $x^{3}-x^{2}-1$   & $1.4656$ & $7,8$\\
$x^{2}=x+3$     & $1,2,5,11,\dots$     & $x^{2}-x-3$  & $2.3028$ & $\varnothing$\\
$x^{4}=x^{3}+1$ & $1,2,3,4,5,7,\dots$  & $x^{4}-x^{3}-1$ & $1.3803$ & $\varnothing$\\
\bottomrule
\end{tabular}
\medskip
\caption{Closure and associativity shifts, computed exactly for $n,m,p\le24$ and $s\le8$. All
$\beta$ listed are Pisot numbers except $\beta=2.3028\dots$, whose conjugate is $-1.3028\dots$;
$\beta=1.3803\dots$ is the second smallest Pisot number.}
\label{tab:main}
\end{table}

\begin{table}[ht]
\centering
\small
\begin{tabular}{lll}
\toprule
system & $u_{0},u_{1},u_{2},\dots$ & $G_{40}\cap[0,14]$\\
\midrule
Zeckendorf     & $1,2,3,5,8,\dots$   & $2,3,4,5,6,7,8,9,10,11,12,13,14$\\
Lucas          & $1,3,4,7,11,\dots$  & $2,\ \ \ \,4,5,6,7,8,9,10,11,12,13,14$\\
shifted Fib.   & $1,4,5,9,14,\dots$  & $4,\ \ \ \,6,7,8,9,10,11,12,13,14$\\
Pell canonical & $1,3,7,17,41,\dots$ & $2,3,4,5,6,7,8,9,10,11,12,13,14$\\
Pell natural   & $1,2,5,12,29,\dots$ & $1,2,3,4,5,6,7,8,9,10,11,12,13,14$\\
\bottomrule
\end{tabular}
\medskip
\caption{Associativity shifts over a longer range, computed exactly for $n,m,p\le40$.}
\label{tab:focus}
\end{table}

\begin{proposition}\label{prop:lucas}
For the Lucas system $\Uu=(1,3,4,7,11,18,29,47,76,123,199,\dots)$, that is $u_{k}=L_{k+1}$, one has
$3\notin G(\Uu)$. Explicitly,
\[
  (1\circ_{3}2)\circ_{3}6=268\ne278=1\circ_{3}(2\circ_{3}6).
\]
\end{proposition}

\begin{proof}
The greedy expansions are $K(1)=\{0\}$, $K(2)=\{0,0\}$ and $K(6)=\{2,0,0\}$, since $6=4+1+1$.
Hence $1\circ_{3}2=u_{3}+u_{3}=14$, and $K(14)=\{4,1\}$ because $14=11+3$, so
\[
  (1\circ_{3}2)\circ_{3}6=\bigl(u_{9}+u_{7}+u_{7}\bigr)+\bigl(u_{6}+u_{4}+u_{4}\bigr)
  =(123+47+47)+(29+11+11)=268 .
\]
On the other side $2\circ_{3}6=2\bigl(u_{5}+u_{3}+u_{3}\bigr)=2(18+7+7)=64$, and
$K(64)=\{7,4,2,0,0\}$ because $64=47+11+4+1+1$, so
\[
  1\circ_{3}(2\circ_{3}6)=u_{10}+u_{7}+u_{5}+u_{3}+u_{3}=199+47+18+7+7=278 . \qedhere
\]
\end{proof}

\begin{proposition}\label{prop:others}
Let $\Uu$ be the system with $u_{0}=1$, $u_{1}=4$ and $u_{k}=u_{k-1}+u_{k-2}$, namely
\[
  1,\,4,\,5,\,9,\,14,\,23,\,37,\,60,\,97,\,157,\,254,\,411,\,665,\,1076,\,1741,\,2817,\,4558,\,7375,\dots
\]
Then $5\notin G(\Uu)$, certified by
$(1\circ_{5}8)\circ_{5}8=8017\ne8072=1\circ_{5}(8\circ_{5}8)$. For the Zeckendorf system
$1\notin G(\Uu)$, certified by $(1\circ_{1}4)\circ_{1}4=40\ne41=1\circ_{1}(4\circ_{1}4)$.
\end{proposition}

\begin{proof}
Both are finite computations of the same shape as Proposition \ref{prop:lucas}; they are carried out
in \S3 and \S4 of \texttt{verify.py}. For the first: $K(8)=\{2,0,0,0\}$, $1\circ_{5}8=u_{7}+3u_{5}=129$,
$K(129)=\{8,5,3\}$, and $(1\circ_{5}8)\circ_{5}8=6045+1427+545=8017$; while
$8\circ_{5}8=u_{9}+6u_{7}+9u_{5}=724$, $K(724)=\{12,6,4,2,0,0,0\}$, and
$1\circ_{5}724=u_{17}+u_{11}+u_{9}+u_{7}+3u_{5}=7375+411+157+60+69=8072$.
\end{proof}

\begin{observation}[$G$ need not be upward closed]\label{obs:notupset}
By Table \ref{tab:focus} and Proposition \ref{prop:lucas}, the Lucas system satisfies
$2\in G$, $3\notin G$, $\{4,\dots,14\}\subseteq G_{40}$; the shifted Fibonacci system satisfies
$4\in G_{40}$, $5\notin G$, $\{6,\dots,14\}\subseteq G_{40}$.
\end{observation}

Lemma \ref{lem:multiples} explains part of this: once $2\in C$, all even shifts lie in $C$. What
fails is the hypothesis of Remark \ref{rem:canonical}. Indeed, for the Lucas system
$2=\Phi(2)\in S$ while $2\varphi\notin S$: here $L(2\varphi)=2u_{1}=6$ and
$\Phi(6)=\varphi^{2}+2\ne2\varphi$, since $K(6)=\{2,0,0\}$. Hence $\varphi S\not\subseteq S$ and
there is no mechanism forcing $3\in C$ --- and in fact $3\notin G$. The odd shifts $5,7,9,\dots$
nevertheless appear to be admissible, and we do not know why; see Question \ref{q:lucasodd}.

\begin{observation}[dependence on initial values]\label{obs:initial}
The three systems with $\beta=\varphi$ in Table \ref{tab:main} have distinct admissible-shift sets.
Hence $G(\Uu)$ is an invariant of the pair $(\beta,\lambda)$ of Proposition \ref{prop:trace}, not of
$\beta$ alone.
\end{observation}

\begin{observation}[non-canonical can beat canonical]\label{obs:beat}
For $\beta=1+\sqrt2$ the system $1,2,5,12,29,\dots$ admits $s=1$, whereas the canonical system
$1,3,7,17,41,\dots$ has least admissible shift $2$. Equivalently: a suitable choice of $\lambda$
produces a set $S$ of $\Uu$-integers closed under multiplication by a smaller power of $\beta$ than
$\Z_{\beta}^{\ge0}$ is.
\end{observation}

\begin{observation}[two ways to fail for all $s$]\label{obs:fail}
For $\beta=2.3028\dots$, the root of $x^{2}=x+3$, no shift $s\le20$ is admissible (checked for
$n,m,p\le12$); $\beta$ is not Pisot, its conjugate being $-1.3028\dots$. Being Pisot is however not
sufficient: for $\beta=1.3803\dots$, the root of $x^{4}=x^{3}+1$ and the second smallest Pisot
number, no shift $s\le20$ is admissible either. This is consistent with the finiteness property
(F) of \cite{FrougnySolomyak}: the digit string of $\beta$ here is $1001$, which violates the
descending-digit criterion, whereas all systems in Table \ref{tab:main} with $C\ne\varnothing$ have
$\beta$ Pisot.
\end{observation}

\begin{observation}[no accidental associativity]\label{obs:noaccident}
In all thirteen systems of Table \ref{tab:main} and all shifts $s\le8$, one has
$C_{24}(\Uu)=G_{24}(\Uu)$.
\end{observation}

\section{Questions}

Theorem \ref{thm:criterion}(iii) proves $C(\Uu)\subseteq G(\Uu)$. Observation
\ref{obs:noaccident} suggests that nothing more is possible.

\begin{conjecture}\label{conj:noaccident}
$C(\Uu)=G(\Uu)$ for every linear numeration system: every associative Knuth product is a transported
ring product. Equivalently, by Theorem \ref{thm:criterion}(iv), the bilinear obstruction
\eqref{eq:obstruction} cannot vanish identically unless the defect $\delta_{s}$ does.
\end{conjecture}

\begin{question}\label{q:lucasodd}
Is $G(\Uu)=\{2\}\cup[4,\infty)$ for the Lucas system, and $\{4\}\cup[6,\infty)$ for
$1,4,5,9,14,\dots$? More pointedly: $3$ fails and $5$ apparently does not, although neither is a
multiple of the least admissible shift. What distinguishes them?
\end{question}

\begin{question}\label{q:zero}
For which $(\beta,\Uu)$ is $0\in C(\Uu)$, i.e.\ for which numeration systems is the set $S$ of
$\Uu$-integers closed under multiplication? Base $b$ (Example \ref{ex:base}) is such a system. Is
$\beta\in\Z$ necessary?
\end{question}

\begin{question}\label{q:invariant}
For fixed $\beta$, put $s_{\min}(\beta)=\min_{\Uu}\min C(\Uu)$, the minimum over all admissible
initial values. Observation \ref{obs:beat} shows $s_{\min}(1+\sqrt2)\le1<L^{\otimes}(1+\sqrt2)$.
Compute $s_{\min}$; it is an invariant of $\beta$ refining $L^{\otimes}$.
\end{question}

\appendix

\section{Verification}\label{sec:verification}

The accompanying script \texttt{verify.py} (Python 3, standard library only) reproduces every finite
assertion above. All arithmetic is exact: integers are Python integers, and elements of
$\Z[\beta]\cong\Z[x]/(\chi)$ are represented by their coordinate vectors in the basis
$1,\beta,\dots,\beta^{d-1}$, with multiplication performed by polynomial multiplication followed by
reduction using $\chi$. No floating-point arithmetic occurs anywhere; the decimal values of $\beta$
quoted in Table \ref{tab:main} are for orientation only and play no role in any check. Membership
$z\in S$ is decided exactly, by testing whether $\Phi(L(z))=z$, which is legitimate by
Corollary \ref{cor:bijection}.

The script performs, in order: verification that each $\chi$ used is irreducible (by exhaustion over
rational roots, and for $d=4$ over factorisations into monic integer quadratics); verification of
Theorem \ref{thm:transport}(i) and (ii) for $n,m\le30$ and $s\le6$ in all thirteen systems;
verification that $C_{24}$ and $G_{24}$ agree, and that closure implies associativity, for $s\le8$
(Table \ref{tab:main}); the extended computation of Table \ref{tab:focus}; the minimal
counterexamples of Propositions \ref{prop:lucas} and \ref{prop:others}, obtained by exhaustive
search in lexicographic order over $n,m,p<60$; Knuth's theorem as a control, for $n,m,p\le45$; the
long-range scans of Observation \ref{obs:fail} for $s\le20$; Lemma \ref{lem:multiples}; the
assertions about $2\varphi\notin S$ used after Observation \ref{obs:notupset}; and
Proposition \ref{prop:trace} for the two systems named there. The script terminates with
\texttt{ALL CHECKS PASSED} and reports no failures.

\begin{thebibliography}{99}

\bibitem{AFMP} P.~Ambro\v{z}, C.~Frougny, Z.~Mas\'akov\'a and E.~Pelantov\'a,
  \emph{Arithmetics on number systems with irrational bases},
  Bull. Belg. Math. Soc. Simon Stevin \textbf{10} (2003), no.~5.

\bibitem{Arnoux} P.~Arnoux, \emph{Some remarks about Fibonacci multiplication},
  Appl. Math. Lett. \textbf{2} (1989), 319--320.

\bibitem{Bertrand} A.~Bertrand, \emph{D\'eveloppements en base de Pisot et r\'epartition modulo~$1$},
  C. R. Acad. Sci. Paris S\'er. A--B \textbf{285} (1977), A419--A421.

\bibitem{BFGK} \v{C}.~Burd\'ik, C.~Frougny, J.-P.~Gazeau and R.~Krejcar,
  \emph{Beta-integers as natural counting systems for quasicrystals},
  J. Phys. A \textbf{31} (1998), 6449--6472.

\bibitem{FrougnySolomyak} C.~Frougny and B.~Solomyak, \emph{Finite beta-expansions},
  Ergodic Theory Dynam. Systems \textbf{12} (1992), 713--723.

\bibitem{Grabner} P.~J.~Grabner et al., \emph{Associativity of recurrence multiplication},
  Appl. Math. Lett. \textbf{7} (1994), 85--90.\footnote{The bibliographic data for this item are
  taken from secondary indexing; we were unable to consult the original, and the author list may be
  incomplete. Its content, as indexed, is the extension of \eqref{eq:knuth} to recurrences
  $G_{k+d}=a_{1}G_{k+d-1}+\dots+a_{d}G_{k}$ with canonical initial values, with associativity proved
  once the index shift is chosen large enough --- that is, the canonical case of
  Theorem \ref{thm:criterion}(iii) together with the fact, our Remark \ref{rem:canonical}, that
  there the admissible shifts form an interval.}

\bibitem{Knuth} D.~E.~Knuth, \emph{Fibonacci multiplication},
  Appl. Math. Lett. \textbf{1} (1988), 57--60.

\bibitem{Messaoudi1} A.~Messaoudi, \emph{G\'en\'eralisation de la multiplication de Fibonacci},
  Math. Slovaca \textbf{50} (2000), no.~2, 135--148.

\bibitem{Messaoudi2} A.~Messaoudi, \emph{Tribonacci multiplication},
  Appl. Math. Lett. \textbf{15} (2002), 981--985.

\bibitem{Parry} W.~Parry, \emph{On the $\beta$-expansions of real numbers},
  Acta Math. Acad. Sci. Hungar. \textbf{11} (1960), 401--416.

\bibitem{Zeckendorf} E.~Zeckendorf, \emph{Repr\'esentation des nombres naturels par une somme de
  nombres de Fibonacci ou de nombres de Lucas},
  Bull. Soc. Roy. Sci. Li\`ege \textbf{41} (1972), 179--182.

\end{thebibliography}

\end{document}
